%%%%%%%%%%%%%%%%%%%%%%%%%%%%%%%%%%%%%%%%%%%%%%%%%%%%%
%			HLAVIČKA								%
%%%%%%%%%%%%%%%%%%%%%%%%%%%%%%%%%%%%%%%%%%%%%%%%%%%%%
\documentclass[openany,12pt]{memoir}

\usepackage[utf8]{inputenc} 
\usepackage[czech]{babel}
\usepackage[T1]{fontenc}
\usepackage[top=1.5cm, bottom=2cm, left=2cm, right=2cm]{geometry}  % --> NASTAVENÍ OKRAJŮ
\usepackage{fancyhdr}
\usepackage{graphicx}
%\usepackage{lmodern}
\usepackage{tgschola}  % --- Silnější font
\usepackage{xwatermark}
\usepackage{xcolor}
\usepackage{changepage}
\usepackage{newunicodechar}



%%%%%% Package na zpěvník
\usepackage[full]{leadsheets}%http://mirrors.nic.cz/tex-archive/macros/latex/contrib/leadsheets/leadsheets_en.pdf   --> dokumentace	
\definesongtitletemplate{empty}{} 
\setchords{
format = \bfseries,   %tučné akordy
minor = {mi},% 
input-notation = {german},%
output-notation = {german}%
}
\definesongtitletemplate{empty}{} 

\newlength{\drop}
\newwatermark[allpages,color=red!50,angle=0,scale=2, xpos=0,ypos=0]{\includegraphics[width=5cm]{obr/pozadi2.jpg}} %--> dvojka na pozadí


%%%% Vlastní příkazy
\newcounter{Slokočet}   %Automatické číslování slok
\newcommand{\mezera}{\vspace*{0.5cm}}   %Horizontální odsazení slok
\newcommand{\stred}{5.2cm}   %%% Na zarovnání slok doprostřed, pozn. automatičtější zarovnávání na střed nejde
\newcommand{\refren}{\mezera \noindent \textbf{R:} } %refrén
\newcommand{\sloka}{\mezera \noindent \addtocounter{Slokočet}{1} \arabic{Slokočet}. } 	%sloka, která se automaticky čísluje
\newcommand{\ssloka}{\mezera \noindent}  % vlastní číslo sloky
\newcommand{\carka}{,\:}
\newcommand{\predehra}{\mezera \hspace*{-1.85cm} \textbf{Předehra:} }
\newcommand{\m}[1]{\color{white}{#1}}  %Pro akordy
\newcommand{\ap}{'} %Pro apostrof
\newcommand{\elipsa}{\kern\fontdimen3\font} %Příkaz pro lepší zacházení s výpustkami (=...); je to vpodstatě jen mezera mezi tečkama výpustky

%%%%%%%%%%%%%%%%%%%%%%%%%%%%%%%%%%%%%%%%%%%%%%%%%%%%%%
%			NÁVOD									 %
%%%%%%%%%%%%%%%%%%%%%%%%%%%%%%%%%%%%%%%%%%%%%%%%%%%%%%
%1. Věci v hlavičce IGNOROVAT
%2. Píseň psát do prostoru mezi \begin{song} a \end{song}
%3. další řádek se značí dvěma odsazeníma (= dvakrát stisknout enter)
%4. \refren vždy na začátku refrenu a \sloka na začátek sloky (automaticky se čísluje)
%  \ = alt gr + q ; [/] = alt gr f/g ; {/} = alt gr + b/n; ^ = alt gr + 3 + mezera
%Cokoliv napíšete do ^{  } se bude brát jako akord
%když se toto bude dotýkat nějakého slova (nebude mezi tím a slovem mezera)
%tak se akordy zjeví nad slovem, ale pište to před slovo
%Když se to nedotýká slova, tak akord lítá ve vzduchu a vytiskne se větší mezera
%První možnost je asi preferovanější
%5. Akordy stačí psát jen do první sloky, když se nezmění -- kytaristi to zvládnou



%\usepackage{subfiles}
%</++>

\begin{document}

\pagestyle{empty}
%\addtocontents{toc}{\protect\thispagestyle{empty}}
%\tableofcontents \thispagestyle{empty}\newpage
\newgeometry{top=1.5cm, bottom = 0cm, left = 2cm, right = 2cm}

\newpage

%Pro kompilaci je důležitý, 	aby ten input a \end{document} zůstali na posledních
%dvou řádcích!!!!
\begin{song}{title=\predtitle\centering Pohár a Kalich \\\large Klíč\vspace*{-0.3cm}}  %% sem se napíše jméno songu a autor
\begin{centerjustified}

\sloka
	^*{Ami}Zvedn ěte poháry, ^*{G}pá nové, už jen ^*{C}po slední ^*{G}př ípitek ^*{Ami}zbývá ,

	pohleďte, nad polem ^*{G}Šp itálským ^*{C}vr cholek ^{G\z}roubení ^*{C}sk rývá.

   	/: Za chvíli ^*{G}zb yde tam popel a ^*{Ami}tráva ,

   	nás čeká ^{G\z }vítězství, bohatství, ^{Ami\z}sláva. :/

\sloka
	Ryšavý panovník jen se smál, sruby prý dobře se boří,

   	palcem k zemi ukázal, ať další hranice hoří.

   	/: Ta malá země už nemá co získat,

	teď bude tancovat, jak budem pískat. :/

\refren
	^{Dmi}Náhle se pozvedl ^*{G}ví tr a mocně ^*{C}vá l,

   	^*{Dmi}odněk ud přinesl ^*{G}ná pěv, sám si ho ^{E7\,\,}hrál, hm\elipsa\dots

\sloka
  	Nad sruby korouhev zavlála, to však neznačí, že je tam sláva,

   	všem věrným pokaždé nesmírnou sílu a jednotu dává.

   	/: Ve znaku kalich, v kalichu víra,

   	jen pravda vítězí, v pravdě je síla. :/

\sloka
	\uv{Modli se, pracuj a poslouchej}, kázali po celá léta,

	Mistr Jan cestu ukázal proti všem úkladům světa.

   	/: Být rovný s rovnými, muž jako žena,

	dávat všem věcem ta pravá jména. :/

\sloka
	Do ticha zazněla přísaha -- ani krok z tohoto místa,

   	se zbraní každý vyčkává, dobře ví, co se chystá.

   	/: Nad nimi stojí muž, přes oko páska,

   	slyšet je dunění a dřevo praská. :/

\refren

   ^{Dmi\,\,}Kdož jsú boží ^*{\z C\,Ami}bojovníci~a~z ákona ^{\z E7\,Ami}jeho,~~~~~

   ^*{Ami}prost ež od ^{\,\,\,\,F}Boha ^{{\color{white}\_\_}Ami}pomoci a ^*{G}úf ajte v ^{\z Ami}něho.~~



\end{centerjustified}
\setcounter{Slokočet}{0}
\end{song}

\end{document}
